% Generated by scripts/cale/paper_numbers.py from aria_pace/results/closed_loop_sensitivity.json (A6). Do not edit by hand.
\begin{table}[t]
\centering
\small
\setlength{\tabcolsep}{3.5pt}
\begin{tabular}{lcccc}
\toprule
Assumed & \multicolumn{3}{c}{Grid points (of 20)} & $\Delta J$ per sequence \\
\cmidrule(lr){2-4}
$p_{\mathrm{stuck}}$ & Quiet & Tie & Always & (Quiet $-$ always) \\
\midrule
0 & 20 & 0 & 0 & $-0.60$ to $-0.52$ \\
0.05 & 20 & 0 & 0 & $-0.29$ to $-0.21$ \\
0.1 & 0 & 5 & 15 & $+0.00$ to $+0.07$ \\
0.2 & 0 & 0 & 20 & $+0.56$ to $+0.62$ \\
0.3 & 0 & 0 & 20 & $+1.07$ to $+1.11$ \\
\bottomrule
\end{tabular}
\caption{Closed-loop simulation (A6; assumed simulated students, not evidence of effect): when does the Quiet policy ($\tau=0.57$) beat always-help? Rows: the assumed probability that a student left unaided after a wrong step abandons, the assumption that drives the sign. Each row pools 5 help-effect values (0--0.5) $\times$ 4 interruption-cost values (0--0.3). Quiet / Always: that policy has the lower cost $J$; tie: the 95\% interval of simulation noise includes 0. Quiet's advantage flips at $p_{\mathrm{stuck}}\approx\AsixSimQuietFlipMin$--\AsixSimQuietFlipMax. Supportive ($\tau=0.09$) makes the same decisions as always-help at every point.}
\label{tab:robust}
\end{table}
